% Template:     Tesis LaTeX
% Documento:    Archivo principal
% Versión:      3.4.8 (02/12/2025)

% CREACIÓN DEL DOCUMENTO
\documentclass[
	spanish, % Idioma: spanish, english, etc.	
	letterpaper, oneside
]{book}

% INFORMACIÓN DEL DOCUMENTO
% DESIGN OF ROBUST CONTROL SYSTEMS
\def\documenttitle {Diseño de sistemas de control robustos}
\def\documentsubtitle {}
\def\degreetitle {
	% FROM CLASSICAL TO MODFRN PRACTICAI. APPROACHES
	Del enfoque clásico al moderno en el diseño de sistemas de control
	\bigbreak\vspace{0.3cm}
	Un enfoque práctico - Marcel J. Sidi
}

\def\universityname {Universidad Autónoma de Nuevo León}
\def\universityfaculty {Facultad de Ingeniería Mecánica y Eléctrica}
\def\universitydepartment {Ingeniería Eléctrica}
\def\universitydepartmentimage {departamentos/fime-logo-cel}
\def\universitydepartmentimagecfg {height=3cm}
\def\universitylocation {San Nicolás de los Garza, Nuevo León, México}

% INTEGRANTES, PROFESORES Y FECHAS
\def\documentauthor {Joanathan de la Cruz}
\def\documentdate {\the\year}

\def\portrait {
	\begin{center}
	\vspace{1.5cm} ~ \\
	\MakeUppercase{\textbf{\documenttitle}} ~ \\
	\vspace{1.5cm}
	\MakeUppercase{\degreetitle} ~ \\
	\vfill
	\begin{tabular}{c}
		\MakeUppercase{\textbf{\documentauthor}} \\ \\
		\vspace{1.0cm} \\
		% PROFESOR GUÍA: \\
		% PROFESOR 1 \\
		% \vspace{0.5cm} \\
		% MIEMBROS DE LA COMISIÓN: \\
		% PROFESOR 2 \\
		% PROFESOR 3 \\
		% \vspace{0.5cm} \\
		% Este trabajo ha sido parcialmente financiado por: \\
		% NOMBRE INSTITUCIÓN \\
		\vspace{0.5cm} \\
		\MakeUppercase{\universitylocation} \\
		\MakeUppercase{\documentdate}
	\end{tabular}
	\end{center}
}
\def\abstracttable {
	% \begin{tabular}{l}
	% 	RESUMEN DE LA MEMORIA PARA OPTAR \\
	% 	AL TÍTULO DE MAGÍSTER EN CIENCIAS \\
	% 	DE LA INGENIERÍA \\
	% 	POR: \MakeUppercase{\documentauthor} \\
	% 	FECHA: \MakeUppercase{\documentdate} \\
	% 	PROF. GUÍA: NOMBRE PROFESOR
	% \end{tabular}
}

% IMPORTACIÓN DEL TEMPLATE
\input{template}

% INICIO DE LAS PÁGINAS
\begin{document}

% PORTADA
\templatePortrait

% CONFIGURACIÓN DE PÁGINA Y ENCABEZADOS
\templatePagecfg

% RESUMEN O ABSTRACT
\begin{abstractd}
	Una de las principales motivaciones para escribir este libro es el creciente abandono de la teoría de diseño clásica frente al tremendo desarrollo de la teoría de control en los últimos veinte años. Con demasiada frecuencia, el estudiante de hoy omite los enfoques de diseño clásico y se enfoca directamente en teorías avanzadas modernas. Esto es desafortunado. Una comprensión profunda de la teoría fundamental clásica de retroalimentación es esencial para entender procedimientos de diseño avanzados como el control óptimo $H_{\infty}$. Para remediar esta deficiencia, nuestro enfoque es introducir progresivamente al lector en técnicas modernas de diseño de control a través de la presentación de teoría de control clásica más tradicional.

	La segunda motivación para escribir el libro es de importancia práctica. Los trabajos avanzados en teoría de control se concentran en desarrollos teóricos y al leer tales libros, uno se fascina por la exposición del autor, las pruebas matemáticas, etc. Sin embargo, con demasiada frecuencia, según la experiencia de este autor con estudiantes, la esencia del material estudiado nunca se comprende completamente, a menos que el estudiante se enfrente a problemas ilustrativos y su solución. Entonces la aplicabilidad de los teoremas enunciados y comprobados se hace aparente. En Diseño de Sistemas de Control Robusto, se atribuye menor importancia a las pruebas de teoremas (estas se pueden encontrar en muchos libros excelentes sobre teoría de control). Por otro lado, cada tema recién tratado es seguido por un ejemplo que ilustra completamente los problemas de diseño, permitiendo al lector repetir la solución y adquirir la experiencia necesaria y comprensión de la técnica presentada.

	El libro enfatiza los aspectos prácticos en el diseño de sistemas de control de retroalimentación en los que la planta puede ser no mínima, inestable e incluso altamente incierta. La anticipación temprana e interpretación correcta de estos efectos puede acortar considerablemente el proceso de diseño. La propiedad de incertidumbre de la planta está estrechamente relacionada con lo que hoy se llama `robustez'. Una tarea importante en control automático es diseñar sistemas controlados de retroalimentación bien comportados que permanezcan `robustos' y satisfagan desempeños cuantitativos formalmente definidos para toda la región de incertidumbre de la planta. Los enfoques de diseño `clásico' y `moderno' para plantas inciertas se explican lado a lado y se utilizan para resolver ejemplos de diseño similares. Esto permite al lector comparar resultados obtenidos por ambos enfoques. De aquí proviene el nombre del libro.

	Los enfoques de diseño presentados en el texto enfatizan fuertemente el principio de `diseño a especificaciones cuantitativas'. Este enfoque elimina parcialmente la necesidad de múltiples iteraciones de diseño de prueba y error hasta lograr una solución satisfactoria.

	El enfoque fundamental de diseño `clásico' presentado aquí, que ha sido desarrollado desde principios de los años setenta, es un área de investigación vinculada a necesidades industriales. También enfatiza la importancia de diseñar con capacidad robusta exactamente para la cantidad de incertidumbre de la planta, pero no más, ahorrando así en ancho de banda y costo de retroalimentación. La técnica de diseño se realiza en el dominio de la frecuencia, que es muy transparente en compensaciones entre diferentes requisitos del sistema. Afortunadamente, también existen buenas aproximaciones para traducir especificaciones del dominio del tiempo al dominio de la frecuencia.

	Al lector se le instruye en cómo resolver eficientemente problemas de control mediante cálculos manuales. Aunque excelente software económico puede usarse hoy en día para propósitos de diseño de control, su uso eficiente depende fuertemente de experiencia práctica de diseño mediante cálculos manuales. Por ejemplo, la caja de herramientas QFT-Matlab se utiliza para resolver el problema de retroalimentación de plantas altamente inciertas en el dominio de la frecuencia. Sin embargo, este software nunca puede utilizarse eficientemente a menos que el diseñador tenga experiencia resolviendo tales problemas mediante cálculos manuales, dibujando manualmente límites permisibles en la función de transmisión de lazo abierto nominal en el plano complejo logarítmico, etc.

	De especial interés son las limitaciones de ganancia-ancho de banda cuando la planta a ser controlada no es de fase mínima, sino que contiene polos inestables y/o ceros en el semiplano derecho. Es importante que los estudiantes se familiaricen con estas limitaciones teóricas de ancho de banda antes de intentar un diseño detallado.

	La teoría de espacio de estados y técnicas de diseño que surgieron a finales de los años cincuenta se formulan en el dominio del tiempo, pero también están estrechamente vinculadas al dominio de la frecuencia y así nos permiten mezclar requisitos y técnicas de diseño en ambos dominios. Técnicas de diseño como el LQR (Regulador Cuadrático Lineal) y LQG (Gaussiano Cuadrático Lineal) han estado en uso común durante mucho tiempo. Se presentan en el texto con algunos ejemplos ilustrativos completamente resueltos. La deficiencia de estas celebradas técnicas de diseño radica en su incapacidad para lidiar con la incertidumbre de plantas restringidas. Esta deficiencia motivó el uso del paradigma de optimización de norma $H_{\infty}$ para resolver el problema de control de plantas inciertas. El enfoque de control $H_{\infty}$, que emergió a principios de los años ochenta, se está difundiendo rápidamente hoy en día entre la comunidad de control. Las características fundamentales de estos nuevos enfoques de diseño también se tratan en el libro. Problemas de control básicos resueltos con técnicas clásicas en el dominio de la frecuencia también se resuelven con técnicas de diseño de optimización de norma $H_{\infty}$. Nuestro deseo es darle al lector la oportunidad de comprender las características comunes de ambas filosofías de diseño, así como apreciar la superioridad, si la hay, de cada enfoque con respecto a este detalle u otro. En la opinión del autor, es lógico ver todas estas técnicas de diseño como complementarias. Una solución correcta y apropiada de un sistema de control de retroalimentación con el enfoque de diseño de control $H_{\infty}$ requiere conocimiento de teoría de control clásica en el dominio de la frecuencia. Para ilustrar este punto, el conocimiento de cómo traducir requisitos técnicos del dominio del tiempo al dominio de la frecuencia y cómo elegir óptimamente las funciones de transferencia ponderadas es de suma importancia si el enfoque de diseño de control $H_{\infty}$ es para producir la mejor solución `práctica'. Además, las deficiencias en una técnica pueden compensarse con características superiores de la otra. Estos puntos importantes se subrayan en el libro mediante los problemas presentados y sus soluciones detalladas. Por otra parte, se intenta en el texto usar una técnica de diseño combinada QFT/$H_{\infty}$ para sistemas de retroalimentación altamente inciertos de dos grados de libertad, tanto SISO como MIMO.

	El libro es el resultado de mi experiencia docente en la Universidad de Tel-Aviv, así como de mi experiencia práctica como ingeniero de control de desarrollo en Israel Aircraft Industries Ltd., donde durante más de 25 años diseñé sistemas de control de vuelo de retroalimentación para aviones, misiles y naves espaciales.

	Un requisito previo para leer este libro es un primer curso en control automático, de tiempo continuo y de datos muestreados, e introducción a sistemas lineales y álgebra lineal. El libro está destinado a estudiantes avanzados de pregrado, estudiantes de posgrado, así como para ingenieros en ejercicio.

	La mayoría de los ejemplos fueron resueltos por software producido por MATHWORKS Inc., incluyendo programas MATLAB relacionados con control, tales como Control System, Robust Control, QFT, $\mu$-Analysis and Synthesis y LMI Control Toolboxes.
\end{abstractd}

% DEDICATORIA
\begin{dedicatory}
	% A study of the practical aspects in designing feedback control systems in which the plant may be non-minimum phase, unstable and also highly uncertain. Classical (QFT) and modern (Hoo) design approaches are explained side-by-side and are used to solve design examples.
	Un estudio de los aspectos prácticos en el diseño de sistemas de control en los que la planta puede ser no mínima, inestable y también altamente incierta. Los enfoques de diseño clásico (QFT) y moderno (H∞) se explican lado a lado y se utilizan para resolver ejemplos de diseño.
	% Una frase de dedicatoria, \\
	% pueden ser dos líneas. \\
	% ~ \\
	% \textbf{Saludos

	\textbf{Sapere aude}
\end{dedicatory}

% AGRADECIMIENTOS
% \begin{acknowledgments}
% 	\lipsum[1]
% \end{acknowledgments}

% TABLA DE CONTENIDOS - ÍNDICE
\templateIndex

% CONFIGURACIONES FINALES
\templateFinalcfg

% ======================= INICIO DEL DOCUMENTO =======================


\chapter{Introducción}

La primera investigación analítica conocida sobre estabilidad fue realizada por Maxwell en 1868, cuando analizó la estabilidad del regulador centrífugo de James Watt para controlar la velocidad, desarrollado en 1788. La teoría del control automático, tal como hoy se entiende, comenzó a tomar forma a inicios de la década de 1930, cuando los amplificadores con retroalimentación para comunicaciones empezaron a usarse comercialmente. El uso de la \textbf{retroalimentación} surgió con el objetivo de hacer que las características del amplificador fueran menos sensibles a las incertidumbres y anomalías físicas de los componentes electrónicos comunes de aquella época. La retroalimentación también permitió conferir propiedades definidas al “amplificador” como un todo.

Como hoy es bien sabido, la retroalimentación también crea serios problemas de estabilidad. El primer trabajo que trató la inestabilidad fue publicado por \textbf{Nyquist (1932)}, quien formuló el célebre criterio de estabilidad de Nyquist. Su trabajo fue ampliado por \textbf{Black (1934)}, que investigó y analizó la estabilidad de los amplificadores con retroalimentación. La teoría del control automático avanzó de manera sustancial gracias al trabajo fundamental de \textbf{Bode (1945)}.

Durante la Segunda Guerra Mundial, la teoría de la retroalimentación se orientó a resolver problemas de control relacionados con equipos bélicos, como la entonces novedosa instrumentación de radar, el control de submarinos y otros sistemas similares. La teoría de control avanzó con rapidez y, hacia finales de la década de 1950, lo que hoy se denomina \textbf{teoría clásica de control} permitía diseñar numerosos sistemas complejos basados en principios de retroalimentación. Más adelante, con la aparición de hardware muy sofisticado —como aeronaves avanzadas, misiles, naves espaciales, plantas químicas complejas y centrales nucleares—, los teóricos del control tuvieron que buscar y desarrollar nuevos enfoques de diseño, basados en lo que hoy llamamos \textbf{teoría moderna de control}, la cual puede entenderse, en cierto sentido, como una continuación natural de la teoría clásica. Por ello, una buena comprensión de la teoría clásica es requisito previo para entender las complejidades inherentes a las técnicas modernas de diseño de control.

Tanto la teoría clásica como la moderna partieron inicialmente del estudio de procesos con parámetros conocidos y fijos. Desafortunadamente, las plantas complejas suelen ser \textbf{inciertas}, en el sentido de que sus parámetros no se conocen con exactitud; solo se conocen los rangos de sus valores. El sistema de control por retroalimentación debe comportarse de acuerdo con desempeños cuantitativamente especificados; en otras palabras, debe comportarse de manera \textbf{robusta} a pesar de las incertidumbres de la planta.

El objetivo de este libro es aplicar enfoques prácticos de diseño, tanto clásicos como modernos, a sistemas de control por retroalimentación inciertos. El material del libro trata únicamente sistemas inciertos lineales e invariantes en el tiempo. Sin embargo, para fines de ingeniería, también resulta aplicable a plantas con parámetros que varían lentamente.


\section{La Estructura y Componentes de un Sistema de Control Realimentado}
Un sistema de control por retroalimentación es una combinación de componentes físicos que, en conjunto, cumplen una misión definida. En la \textbf{Figura 1.1} se muestra una estructura generalizada de un sistema de control por retroalimentación (SISO o MIMO), seguida de la explicación de sus distintos componentes.

\insertimage[\label{img:estructura-general-retroalimentacion}]{2024_06_29_804c621bf2d98037ffbeg-015_451_1134_1462_737}{width=\textwidth}{Estructura general de retroalimentación de dos grados de libertad. El elemento \(G\) se denomina a veces red de control, compensador de control o simplemente controlador/compensador.}

\subsection{Los componentes de un sistema de control realimentado}
Se distinguen, básicamente, dos tipos de componentes funcionales:

\begin{enumerate}
  \item \textbf{Los componentes pertenecientes a la planta física definida que debe controlarse}, de modo que se alcancen los objetivos especificados del sistema. La \textbf{planta} denota la parte condicionada del sistema, cuyas salidas son las salidas del sistema. Puede describirse mediante un solo bloque de función de transferencia o, si es necesario y factible, mediante varios bloques de subplantas, \(P_1, P_2, \ldots, P_n\), de modo que puedan obtenerse beneficios adicionales de control al medir sus salidas.

  \item \textbf{Los componentes pertenecientes al controlador del sistema de retroalimentación}, que son inherentes y una parte indispensable del sistema controlado, representados por los bloques de función de transferencia \(S, H, G, F_f, F_c\).

\end{enumerate}

\textbf{Descripción de los bloques principales:}

\begin{itemize}
  \item \textbf{\(G\)}: compensador dinámico cuya tarea principal es proporcionar al sistema de retroalimentación la estabilidad dinámica necesaria y los coeficientes definidos de error en estado estacionario. El compensador dinámico también se utiliza para satisfacer especificaciones definidas de lazo cerrado. Es un componente esencial de control en cualquier sistema con retroalimentación.
  \item \textbf{\(S\)}: hardware del sensor que mide las salidas de la planta utilizadas para controlar el sistema de retroalimentación.
  \item \textbf{\(H\)}: filtro dinámico cuyos objetivos principales son procesar las salidas de los sensores. Su función básica es atenuar el ruido inherente de los sensores y las señales asociadas con dinámicas estructurales de la planta (si estas existen), entre otros efectos.
  \item \textbf{\(F_f\)}: prefiltro usado para moldear convenientemente el desempeño de seguimiento entrada-salida del sistema controlado. Si puede implementarse físicamente, proporciona al sistema un \textbf{segundo grado de libertad} para fines de diseño.
  \item \textbf{\(F_c\)}: filtro de precomando que acelera intencionalmente la salida del sistema al acceder directamente a la entrada de la planta.
\end{itemize}

No hace falta decir que estos componentes del controlador pueden implementarse mediante hardware electrónico o mediante software en microprocesador.

\subsection{Las señales en un sistema de control realimentado}
Se distinguen distintas señales básicas, tales como señales externas de referencia y de perturbación, salidas del sistema y señales de control a la entrada de la planta.

\begin{itemize}
  \item \textbf{Entradas de referencia}: pueden ser de dos tipos:\\
(i) entradas de mando generadas por operadores humanos, como un piloto, por ejemplo; o\\
(ii) entradas de mando generadas por la cinemática de un cuerpo en movimiento, como en un sistema de radar que sigue un objeto volador.

  \item \textbf{Entradas de perturbación}: son señales no deseadas que tienden a afectar al sistema controlado, actuando sobre la planta en diferentes puntos de entrada.

  \item \textbf{Señales de ruido}: se producen dentro de los sensores físicos. Pueden causar deterioro en el desempeño del sistema de retroalimentación controlado.

  \item \textbf{Señales de salida}: son las variables que deben controlarse. En un sistema MIMO, la salida es un vector cuyos elementos son las variables dinámicas de la planta. En general, no todas ellas se miden ni se especifican como controlables; incluso algunas podrían no ser medibles.

  \item \textbf{Señales de control en la entrada de la planta}: a veces llamadas \textbf{esfuerzo de control}, son generadas por el compensador \(G\) e inyectadas a la entrada de la planta. Su nivel de amplitud debe mantenerse tan bajo como sea posible para evitar la saturación de la entrada de la planta.

\end{itemize}

\noindent\rule{\textwidth}{0.5pt}

\section{¿Por qué Realimentación?}

¿Por qué necesitamos aplicar retroalimentación para controlar un sistema físico cuya planta es matemáticamente conocida? Para responder de forma sencilla a esta pregunta, se simplifica la estructura de retroalimentación de la Figura 1.1 a la de la \textbf{Figura 1.2a}, donde \(P(s)\) se supone una planta SISO cuya salida es \(y(s)\).

\begin{images}[\label{img:figura-1-2}]{Obtención de características entrada-salida sin retroalimentación. La idea central de la figura es mostrar que, en ausencia de perturbaciones e incertidumbre, podría intentarse imponer una dinámica deseada mediante un bloque directo delante de la planta.}
  \addimage{2024_06_29_804c621bf2d98037ffbeg-016_272_755_145_1069}{width=0.45\textwidth}{a.}
  \imagesnewline
  \addimage{2024_06_29_804c621bf2d98037ffbeg-016_285_1134_469_737}{width=0.45\textwidth}{b.}
\end{images}

Supóngase que la meta es obtener, para el sistema controlado, una función de transferencia entrada-salida especificada:

\[
T(s) = \frac{y(s)}{r(s)} \tag{1.2}
\]

Si la perturbación \(d(s)=0\), esta meta puede alcanzarse agregando delante de la planta una red de control

\[
G(s)=\frac{T(s)}{P(s)}
\]

como en la Figura 1.2b. Esto equivale a cancelar la dinámica de la planta y sustituirla por la dinámica especificada de \(T(s)\).

\subsection{¿Qué tiene de malo este enfoque?}
Aunque existen muchos argumentos a favor del uso de retroalimentación, aquí se presentan los dos principales:

\subsection{A. Sensibilidad a incertidumbre en la planta}
Para una planta perfectamente conocida, este procedimiento podría resultar satisfactorio. Pero, si la planta no se conoce exactamente, entonces \(G(s)\) será una solución errónea. Por tanto, la respuesta entrada-salida no será lo suficientemente parecida a la deseada. En casos extremos, si la planta \(P(s)\) incluye un polo inestable ubicado en \(+a\), incluso una incertidumbre muy pequeña \(\varepsilon\) en \(+a\) hará que la función de transferencia entrada-salida contenga un dipolo inestable de la forma

\[
\frac{s-a+\varepsilon}{s-a}
\]

porque ya no podrá lograrse una cancelación perfecta del polo inestable. En consecuencia, la salida de la planta divergerá inevitablemente.

\subsection{B. Imposibilidad práctica de cancelar perturbaciones desconocidas}
Si la perturbación \(d(s)\) fuera conocida, su efecto en la salida de la planta podría anularse inyectando a la entrada de la planta la señal corregida siguiente:

\[
 u_p(s) = \frac{r(s)T(s)}{P(s)} - d(s), \qquad y(s)=P(s)u_p(s) \tag{1.3}
\]

Entonces la salida sería la respuesta deseada \(y(s)=T(s)r(s)\). Pero, desafortunadamente, la perturbación \(d(s)\) rara vez es bien conocida y sus características por lo general no son deterministas. Por ello, la salida real puede desviarse fuertemente de la respuesta deseada debida únicamente a la entrada de mando \(r(s)\).

\noindent\rule{\textwidth}{0.5pt}

\section{Elementos a Considerar por Diseñadores de Sistemas de Control Realimentado}
La estructura generalizada de retroalimentación se muestra en la Figura 1.1. Las ventajas obtenidas mediante retroalimentación son invaluables. Sin embargo, la propia retroalimentación introduce problemas prácticos. Supóngase que debemos usar retroalimentación alrededor de una planta estable para lograr ciertas características deseadas del sistema controlado. Si el sistema está mal diseñado, puede volverse inestable, aun cuando la planta sea, en esencia, un sistema físico estable. Comprender la estructura de retroalimentación ayuda a entender cuáles son los problemas de diseño y cómo serán tratados en este libro.

\subsection{Especificaciones de diseño}
La realimentación se usa para lograr características y desempeños especificados que, como ya se explicó, no pueden alcanzarse en lazo abierto. Principalmente se busca:

\begin{enumerate}
  \item \textbf{especificaciones cuantitativas de seguimiento entrada-salida}, y
  \item \textbf{atenuación del efecto de perturbaciones externas esperadas} (aunque no exactamente conocidas), ya sea en la salida o en la entrada de la planta.
\end{enumerate}

Las plantas LTI que no incluyen polos ni ceros en el semiplano derecho (RHP) pueden controlarse mediante un controlador LTI para alcanzar cualquier especificación deseada de lazo cerrado, siempre que el nivel de ruido del sensor sea despreciable y que la planta tenga la capacidad física de proporcionar el esfuerzo de control requerido.

Los sistemas estabilizados con plantas que contienen \textbf{polos en RHP} (pero no ceros en RHP) pueden sufrir un ancho de banda del lazo abierto demasiado elevado, intolerable desde el punto de vista de las limitaciones físicas. Lo contrario ocurre en sistemas estabilizados cuya planta contiene \textbf{ceros en RHP}.

En resumen, cuando la planta contiene polos y/o ceros en el semiplano derecho, existen \textbf{limitaciones teóricas} sobre las especificaciones de lazo cerrado que pueden alcanzarse. Una tarea importante de este libro es tratar estas limitaciones de manera cuantitativa.

\subsection{Incertidumbre de los parámetros de la planta}
Los parámetros de la planta no se conocen exactamente; es decir, son inciertos y solo se conocen los rangos de sus valores. El desempeño nominal de lazo cerrado no cambiará drásticamente ante variaciones pequeñas de los parámetros de la planta. Sin embargo, si esos parámetros cambian de manera significativa, esto puede generar problemas serios en el diseño de sistemas de control por retroalimentación, en los que deben mantenerse la \textbf{estabilidad} y los \textbf{desempeños especificados}, a pesar de la incertidumbre conocida de la planta.

Este es el llamado \textbf{problema de robustez}, uno de los principales objetivos de análisis y diseño del libro. Se mostrará que, si la planta incierta incluye ceros en RHP, no existe garantía de que las especificaciones inicialmente definidas puedan alcanzarse.

\subsection{Saturación de la planta}
Cuando deben alcanzarse especificaciones exigentes, el esfuerzo de control requerido puede sobrepasar las capacidades de la planta. En ese caso aparecen efectos no lineales que deben tratarse de manera adecuada.

Si la planta puede subdividirse físicamente en dos o más subplantas en cascada, como en la Figura 1.1, entonces existe una ventaja en usar la salida de \(P_2\), si esta es medible, con el fin de obtener un mejor desempeño cerrando un lazo interno de retroalimentación.

\subsection{Amplificación del ruido del sensor}
Los sensores siempre están presentes en un sistema de control realimentado. Desafortunadamente, producen ruido eléctrico interno que se superpone a la señal medida y útil que debe enviarse al controlador \(G(s)\). Este ruido puede tener un efecto enorme en el diseño de cualquier sistema de retroalimentación.

Como se verá, para lograr tanto baja sensibilidad del sistema de lazo cerrado frente a incertidumbres paramétricas como buen rechazo de perturbaciones externas, se necesita una \textbf{alta ganancia de lazo abierto} en un amplio rango de frecuencias.

Por ejemplo, considérese el problema de atenuar la perturbación \(d(s)\) en la salida del sistema \(y(s)\) en la Figura 1.1:

\[
 y(s)=\frac{d(s)}{1+G(s)H(s)P(s)S(s)}=\frac{d(s)}{1+L(s)} \tag{1.4}
\]

Es claro que, para mantener pequeña a \(y(s)\) frente a la perturbación externa \(d(s)\), la magnitud de

\[
L(j\omega)=G(j\omega)H(j\omega)P(j\omega)S(j\omega)
\]

debe ser grande en un rango de frecuencias suficientemente amplio. Como la planta \(P(s)\) está dada y no puede modificarse, queda claro que la magnitud de \(G(j\omega)H(j\omega)S(j\omega)\) debe incrementarse para disminuir el nivel de \(y(s)\), es decir, para lograr una buena atenuación de perturbaciones.

Pero aumentar la magnitud de \(G(j\omega)H(j\omega)S(j\omega)\) tiene una desventaja importante: el ruido del sensor \(n(s)\) se amplifica a la entrada de la planta de acuerdo con

\[
 u(s)=\frac{-G(s)H(s)}{1+G(s)H(s)P(s)S(s)}n(s)= -\frac{G(s)H(s)}{1+L(s)}n(s) \tag{1.5}
\]

Ahora considérese dos rangos de frecuencia.

\subsubsection{(i) Rango donde \(|L(j\omega)| \gg 1\)}
Generalmente esto aplica a bajas frecuencias, donde suele concentrarse la mayor parte de la energía de la perturbación. En ese caso,

\[
 u(s) \cong -\frac{n(s)}{P(s)S(s)}, \qquad |L(j\omega)| \gg 1 \tag{1.6}
\]

La Ec. (1.6) muestra que, en este rango de frecuencia, la amplificación del ruido no depende realmente del controlador diseñado \(G(s)\) ni del filtro del sensor \(H(s)\), sino exclusivamente de las propiedades de la planta \(P(s)\) y del hardware del sensor \(S(s)\).

\subsubsection{(ii) Rango donde \(|L(j\omega)| \ll 1\)}
En ese caso,

\[
 u(s) \cong -G(s)H(s)n(s), \qquad |L(j\omega)| \ll 1 \tag{1.7}
\]

La Ec. (1.7) muestra cualitativamente que, cuanto mayor sea la magnitud de \(G(j\omega)H(j\omega)\) en este rango, mayor será el efecto de la amplificación del ruido del sensor en la entrada de la planta. En general, la energía del ruido del sensor se concentra justamente en esta región. Así, controladores de mayor ancho de banda amplían el rango de frecuencias en el que el efecto del ruido blanco del sensor resulta significativo.

\subsection{Conclusión práctica de esta sección}
Diseñar el compensador de control para lograr un buen rechazo de perturbaciones externas empeora, en sentido contrario, la amplificación del ruido del sensor y todos los problemas asociados con ella. Por tanto, uno de los desafíos más básicos del control por retroalimentación es encontrar controladores satisfactorios que logren un compromiso entre estos problemas contrapuestos, presentes en la mayoría de los sistemas físicos reales. Este es uno de los temas centrales del libro.

En síntesis, las técnicas de diseño con retroalimentación deben considerar cuantitativamente:

\begin{itemize}
  \item la sensibilidad aceptable de la función de transferencia de seguimiento entrada-salida ante la incertidumbre de la planta;
  \item el rechazo de perturbaciones externas;
  \item las limitaciones del sistema debidas a polos y ceros en RHP presentes en la planta; y
  \item la amplificación del ruido del sensor, que, si no se atenúa adecuadamente, puede saturar la entrada de la planta.
\end{itemize}

\noindent\rule{\textwidth}{0.5pt}

\section{Contenido del Libro}
El libro trata sistemas lineales de retroalimentación \textbf{SISO} (una entrada, una salida) y \textbf{MIMO} (múltiples entradas, múltiples salidas). Las propiedades fundamentales de la retroalimentación se explican y comprenden con mayor claridad mediante sistemas SISO. Además, dado que algunas de las técnicas de diseño MIMO que se estudiarán se apoyan en procedimientos de diseño SISO, es natural dedicar los primeros capítulos principalmente a sistemas SISO, aunque se hará referencia a sistemas MIMO cuando sea necesario. Los capítulos restantes están dedicados a sistemas MIMO.

A continuación se presenta una breve descripción del contenido de cada capítulo, procurando resaltar las interrelaciones más importantes entre ellos.

\subsection{Capítulo 2}
Este capítulo comienza definiendo el problema de retroalimentación. Después se tratan e ilustran propiedades básicas de la retroalimentación. A continuación se introduce y analiza la \textbf{función de sensibilidad entrada-salida}, tan importante para la calidad del sistema de control diseñado. Luego se define la \textbf{estabilidad interna} de los sistemas de retroalimentación y se demuestra un teorema básico relacionado con ese tema.

Se revisan los conceptos básicos de diseño. Se explican los principios de \textbf{loop shaping}. Los enfoques de diseño del libro se basan en gran medida en técnicas en el dominio de la frecuencia. Primero se enuncia el criterio fundamental de estabilidad de Nyquist y después se define la \textbf{estabilidad relativa} clásica. Luego se estudian técnicas de diseño orientadas por Nyquist usando diagramas de \textbf{Bode}, la carta de \textbf{Nichols} y la carta de Nichols invertida.

Se explican e ilustran, con ejemplos completamente resueltos, técnicas básicas de loop shaping para lograr márgenes de estabilidad, anchos de banda especificados y atenuación suficiente tanto de perturbaciones externas como del ruido del sensor.

También se discuten estructuras de retroalimentación de \textbf{uno} y de \textbf{dos grados de libertad}, y se ejemplifica detalladamente el loop shaping para ambas. Se enuncian y comparan las limitaciones y superioridades de ambas estructuras.

Se introducen los compromisos de diseño en sistemas con retroalimentación. Después se analizan requisitos de desempeño tales como márgenes de ganancia y de fase, expresados en términos de las funciones de sensibilidad y de transferencia de unidad-retroalimentada. A la definición de las sensibilidades \textbf{ponderada} y \textbf{mixta} le sigue la definición del \textbf{problema estándar del regulador óptimo \(H_{\infty}\)}. Finalmente, se ilustra el loop shaping mediante el uso de un algoritmo de optimización en norma \(H_{\infty}\), con un ejemplo resuelto en detalle.

\subsection{Capítulo 3}
Este capítulo trata características básicas y temas prácticos en el diseño de sistemas SISO con retroalimentación. Las especificaciones de diseño suelen darse en el dominio del tiempo. Sin embargo, dado que las técnicas de diseño se realizan principalmente en el dominio de la frecuencia, primero se traducen a especificaciones en frecuencia las especificaciones prácticas en tiempo para sistemas de fase mínima y no mínima.

Se desarrollan técnicas de traducción del dominio del tiempo al dominio de la frecuencia y los resultados se resumen en gráficas apropiadas y relaciones analíticas aproximadas. Existe una traducción directa y exacta del dominio del tiempo al de la frecuencia si se conoce la parte real de una función de transferencia. Desafortunadamente, por razones prácticas, los ingenieros de control trabajan mayormente con la \textbf{magnitud} de la función de transferencia, para la cual solo existen relaciones aproximadas entre frecuencia y tiempo. Esas aproximaciones se enuncian y analizan en este capítulo.

Se discute el problema de lograr, de forma óptima, un ancho de banda de ganancia especificado para funciones de transferencia de fase mínima, a través de la llamada \textbf{Característica Ideal de Bode}. Los sistemas que incluyen ceros y/o polos en RHP, y/o retardo puro, sufren limitaciones inherentes para alcanzar anchos de banda deseados. Estas limitaciones se analizan y se obtienen resultados analíticos para predecir los mejores anchos de banda alcanzables cuando tales polos y ceros en RHP están incluidos en la función de transferencia de la planta. Se presentan además algunos resultados nuevos, ilustrados con soluciones numéricas detalladas.

La \textbf{estabilización} es un tema que merece especial atención. La \textbf{parametrización de Youla} (o \textbf{Q-parametrización}) es una técnica de síntesis para encontrar todos los controladores estabilizantes para una planta dada. Esta técnica, explicada aquí para sistemas SISO con retroalimentación, se ilustra con ejemplos resueltos.

Cuando la planta incluye varios polos y ceros en RHP, suele ser difícil elegir la estructura matemática correcta del compensador que estabilizará el sistema de retroalimentación. Los procedimientos mecanizados que aparecen en la literatura para enfrentar este problema suelen producir una solución estable, pero en general con ganancia de lazo abierto de ancho de banda muy limitado y con márgenes de fase y ganancia pobres, lo cual implica malas propiedades de retroalimentación. Se explican e ilustran con ejemplos detallados procesos de diseño orientados a obtener las mejores propiedades de retroalimentación posibles a partir de una planta difícil, inestable y no mínima fase.

\subsection{Capítulo 4}
El material de los capítulos previos trata la síntesis y el diseño de sistemas de retroalimentación en los que la planta está bien definida y sus parámetros son fijos y conocidos. En sistemas reales, los parámetros ni se conocen bien ni pueden considerarse completamente invariantes en el tiempo. Cuando la planta lineal es altamente incierta, las técnicas de diseño de los capítulos 2 y 3 ya no pueden usarse con éxito.

En este capítulo se introduce una planta altamente incierta cuyos parámetros no se conocen con precisión. Solo se conocen los rangos de sus valores. Se desarrollan técnicas de diseño para satisfacer las especificaciones a pesar de las incertidumbres de la planta. Aquí aparece formalmente el tema de la \textbf{robustez}. La teoría detrás de esta técnica se conoce hoy como \textbf{Teoría Cuantitativa de la Retroalimentación (QFT)} y se basa en la teoría clásica de control en frecuencia.

La técnica fundamental presentada para plantas inciertas de fase mínima puede emplearse para obtener una solución de ingeniería para cualquier conjunto de especificaciones a priori del sistema controlado, naturalmente al costo de una función de transmisión de lazo con ancho de banda amplio.

Cuando se consideran plantas con ceros en RHP o sistemas de retroalimentación muestreados, la técnica fundamental debe modificarse ligeramente. Además, debido a limitaciones inherentes en sistemas que incluyen ceros en RHP, puede ocurrir incluso que el diseño no logre cumplir todas las especificaciones inicialmente planteadas. Completar el diseño requiere relajar el desempeño inicialmente exigido e intentar una iteración adicional.

La técnica fundamental, con todas sus variantes para sistemas de fase no mínima y sistemas muestreados, se ilustra con ejemplos prácticos resueltos sistemáticamente y evaluados mediante simulaciones en el dominio del tiempo y en el de la frecuencia.

También se trata en este capítulo el diseño de sistemas inciertos basado en el enfoque de control óptimo \(H_{\infty}\). Se introducen y ejemplifican modelos de incertidumbre de planta compatibles con el paradigma de la norma \(H_{\infty}\). Luego se aborda el diseño para estabilidad robusta y para desempeño robusto, también mediante ejemplos detallados. Finalmente, el problema del sistema SISO incierto lineal con dos grados de libertad se resuelve con el enfoque de control óptimo \(H_{\infty}\) usando principios de QFT. Tanto el diseño QFT como el diseño QFT/\(H_{\infty}\) concluyen con soluciones de calidad comparable.

\subsection{Capítulo 5}
Este capítulo trata plantas en cascada de entrada única. La existencia de una variable interna adicional medible puede aprovecharse en el diseño de un sistema de control por retroalimentación \textbf{SIMO} (una entrada, múltiples salidas) para alcanzar el desempeño especificado con un nivel reducido de ruido de sensor amplificado en la entrada de la planta. Cada variable medible adicional permite cerrar un lazo de retroalimentación interno adicional, reduciendo así la carga de control sobre el lazo principal externo. Esta posibilidad es especialmente importante cuando se trabaja con plantas inciertas.

\subsection{Capítulo 6}
Este capítulo trata la síntesis y el diseño de sistemas MIMO LTI con retroalimentación. El control de sistemas MIMO es una de las áreas más investigadas de las últimas décadas. Se aborda tanto en el marco del espacio de estados como en el dominio de la frecuencia y del plano \(s\).

Primero se explican e ilustran algunas técnicas comunes de diseño para sistemas MIMO “ciertos”, en el plano \(s\) y en frecuencia, tales como la diagonalización de la matriz de transferencia de lazo abierto, la no interacción mediante inversión de la matriz de transferencia de la planta y la técnica de cierre secuencial de lazos.

También se enuncian resultados generales sobre \textbf{estabilidad interna} en sistemas MIMO. Luego se explica la \textbf{Q-parametrización} como preludio a la teoría de los algoritmos que resuelven el problema estándar del regulador óptimo \(H_{\infty}\).

El objetivo principal de este capítulo es presentar herramientas cuantitativas de diseño para sintetizar un controlador alrededor de una planta MIMO altamente incierta, con el propósito de lograr especificaciones deseadas de seguimiento entrada-salida. Cuando los sistemas MIMO inciertos se tratan en el marco de QFT, la idea básica es dividir el proceso de diseño en una secuencia de \(n\) diseños SISO, donde \(n\) es el orden del sistema MIMO. La solución al problema original es entonces la combinación de las soluciones de cada etapa secuencial con las técnicas SISO desarrolladas en el Capítulo 4.

\subsection{Capítulo 7}
La finalidad principal del Capítulo 8 es resolver el problema MIMO incierto con dos grados de libertad en el contexto de la optimización en norma \(H_{\infty}\). En este caso, aunque los requisitos de desempeño técnico también se definen en el dominio de la frecuencia, los algoritmos fundamentales de optimización \(H_{\infty}\) tienen sus raíces en la formulación en \textbf{espacio de estados}.

Además, como aquí se enfrenta un problema MIMO, es esencial contar al menos con una comprensión cercana —aunque sea introductoria— de la teoría básica de optimización para diseño de control MIMO, culminando en la solución del problema del \textbf{Regulador Cuadrático Lineal (LQR)}. Eso es precisamente lo que se desarrolla en este capítulo.

Después, las técnicas LQR se adaptan para resolver problemas de seguimiento de trayectorias, como el seguimiento implícito y explícito de modelo. Luego se revisa brevemente el problema \textbf{Lineal-Cuadrático-Gaussiano (LQG)} para aquellos casos en que algunos estados no son medibles. Como esos estados no pueden medirse, deben estimarse de alguna manera óptima para permitir la implementación de la ley de control óptima obtenida con la técnica LQR. La solución de estos problemas se vincula con la solución de ecuaciones matriciales de Riccati. Por lo general, las demostraciones de teoremas importantes se omiten, pues pueden encontrarse en excelentes libros de texto y artículos técnicos referenciados cuando se requiere.

Comprender estos temas es una etapa necesaria antes de abordar la solución del problema MIMO incierto con dos grados de libertad dentro del contexto de optimización en norma \(H_{\infty}\).

\subsection{Capítulo 8}
Este capítulo trata el problema general de retroalimentación MIMO incierta resuelto con el paradigma de optimización en norma \(H_{\infty}\).

Primero se explican algunos hechos elementales sobre \textbf{valores singulares}, así como su uso para medir el tamaño de matrices de transferencia y para definir desempeños en sistemas MIMO. Después se plantea el \textbf{problema estándar del regulador \(H_{\infty}\)} y se bosqueja la matemática de un algoritmo de optimización para resolverlo.

Finalmente se presenta una técnica de diseño para resolver el problema MIMO incierto de dos grados de libertad. El mismo ejemplo de planta MIMO incierta resuelto en el Capítulo 6 con la técnica clásica QFT se resuelve aquí dentro del paradigma de control \(H_{\infty}\), complementado con principios de QFT. Luego se comparan los desempeños alcanzados con ambas técnicas de diseño.

\subsection{Apéndices}
Los apéndices del libro tienen la intención de resumir algunos aspectos importantes y bien conocidos de la teoría de sistemas lineales usados a lo largo del texto. Naturalmente, no incluyen demostraciones rigurosas.

\begin{itemize}
  \item \textbf{Apéndice A}: trata los \textbf{grafos de flujo de señal}. El resultado final es la conocida \textbf{fórmula de ganancia de Mason}, útil para resolver gráficamente ecuaciones lineales simultáneas, calcular transferencias entre señales y obtener sistemáticamente funciones de transferencia de estructuras de retroalimentación complicadas.

  \item \textbf{Apéndice B}: al tratar problemas MIMO, especialmente aquellos resueltos con algoritmos de optimización en norma \(H_{\infty}\), se usan ampliamente nociones de norma vectorial y matricial. Aquí se enuncian y explican igualdades y desigualdades de valores singulares. También se definen nociones básicas de la formulación en espacio de estados, así como eigenvalores y eigenvectores de matrices reales de funciones de transferencia, para su uso en conexión con el criterio generalizado de estabilidad de Nyquist. Además, se definen transformaciones lineales fraccionales y relaciones asociadas utilizadas en la solución del regulador óptimo \(H_{\infty}\).

  \item \textbf{Apéndice C}: cuando el diseño de control se realiza en el dominio de la frecuencia, el compensador se obtiene durante la etapa de loop shaping agregando polos y ceros reales de primer orden, polos y ceros complejos de segundo orden, y también redes lead-lag o lag-lead. Otro recurso muy útil es una red de control compuesta por dos polos y dos ceros colocados de tal manera que se obtienen características especiales de ganancia y fase durante el loop shaping. Para ello se incluyen gráficas y tablas útiles. También se definen la carta de Nichols y su versión invertida, junto con una explicación práctica de su uso.

  \item \textbf{Apéndice D}: presenta, para conveniencia del lector, definiciones y teoremas básicos de las teorías de \textbf{Fourier}, \textbf{Laplace} y \textbf{\(Z\)}.

  \item \textbf{Apéndice E}: la teoría del control por retroalimentación trata con funciones analíticas \(F(s)\) en el plano \(s\). \(F(s)\) puede describir la dinámica de un proceso lineal —usualmente llamado planta— o representar la función de transmisión de lazo abierto o de lazo cerrado en sistemas de retroalimentación. De hecho, todas las teorías de control en el dominio \(s\) se basan en este tipo de funciones.

Cuando se trabaja en el dominio práctico de la frecuencia, a cada valor de \(s=j\omega\) le corresponde un valor de la función en el plano complejo, expresado como
\[
F(j\omega)=A(\omega)+jB(\omega)
\]
Las partes real e imaginaria de estas funciones desempeñan un papel importante en el análisis de sistemas lineales. Además, existen relaciones muy importantes entre esas partes real e imaginaria, con fuerte impacto sobre propiedades prácticas de retroalimentación, como el ancho de banda máximo alcanzable. Este apéndice resume relaciones útiles, de carácter práctico, entre las partes real e imaginaria de funciones analíticas usadas en teoría de control, y también señala resultados importantes sobre restricciones teóricas de dichas funciones. El material sigue de cerca el trabajo fundamental de Bode (1945), así como aportaciones de Guillemin (1949), Seshu y Balabanian (1959) y Horowitz (1963).

  \item \textbf{Apéndice F}: investiga cuál es el \textbf{orden mínimo de un compensador} necesario para lograr la estabilización de un sistema de retroalimentación específico cuando se utiliza un procedimiento de ubicación de polos, como en una técnica de diseño presentada en el Capítulo 3.

  \item \textbf{Apéndice G}: los \textbf{coeficientes de error en estado estacionario} desempeñan un papel importante al preparar especificaciones para un sistema controlado por retroalimentación. Aquí se definen e ilustran los coeficientes usuales para entradas de posición, velocidad y aceleración en distintos tipos de sistemas de retroalimentación.

\end{itemize}


\chapter{Propiedades básicas y diseño de sistemas de retroalimentación SISO}

Los \textbf{Capítulos 2 a 5} del libro se concentran en sistemas de retroalimentación \textbf{SISO}. El autor parte de la idea de que muchas propiedades fundamentales de la retroalimentación se comprenden mejor en este contexto, y que después pueden extenderse al caso \textbf{MIMO}. En este capítulo se fijan las definiciones básicas, se presenta la formulación matemática de la retroalimentación, se introducen las herramientas clásicas de diseño en frecuencia (\textbf{Nyquist, Bode y Nichols}) y, finalmente, se conecta ese enfoque con la formulación moderna basada en \textbf{optimización con norma \(H_\infty\)}.


\section{Definiciones introductorias}
La estructura de retroalimentación que se utilizará a lo largo del capítulo es una forma canónica de \textbf{dos grados de libertad (TDOF)} con retroalimentación unitaria. En comparación con la estructura general del capítulo anterior:

\begin{itemize}
  \item el bloque \(H(s)\) se omite porque, desde el punto de vista de las propiedades de retroalimentación, puede incorporarse a \(G(s)\);
  \item el filtro de precomando \(F_c(s)\) se omite por simplicidad, ya que no altera las propiedades fundamentales del lazo cerrado;
  \item las subplantas en cascada \(P_1\) y \(P_2\) se agrupan en una sola planta \(P(s)\).
\end{itemize}

En esta estructura, la entrada de mando \(r(t)\) y la salida \(y(t)\) pueden medirse de manera independiente. Bajo esa hipótesis, dos compensadores de diseño —\(G(s)\) y \(F(s)\)— permiten actuar sobre dos especificaciones distintas del sistema.

\insertimage[\label{img:estructura-canonica-retroalimentacion}]{2024_06_29_804c621bf2d98037ffbeg-021_337_1135_202_763}{width=0.75\textwidth}{Estructura canónica de retroalimentación unitaria de dos grados de libertad.}

Las funciones de transferencia relevantes para esta estructura son:

\[
T(s)\stackrel{\text{def}}{=}\frac{y(s)}{r(s)}=F(s)\frac{L(s)}{1+L(s)}, 
\qquad 
L(s)\stackrel{\text{def}}{=}G(s)P(s) \tag{2.1.1}
\]

\[
T_{d1}(s)\stackrel{\text{def}}{=}\frac{y(s)}{d_1(s)}=\frac{P(s)}{1+L(s)} \tag{2.1.2}
\]

\[
T_d(s)\stackrel{\text{def}}{=}\frac{y(s)}{d(s)}=\frac{1}{1+L(s)}\equiv S(s) \tag{2.1.3}
\]

\[
T_{un}(s)\stackrel{\text{def}}{=}\frac{u(s)}{n(s)}
=\frac{-G(s)}{1+L(s)}
=-G(s)S(s)
=-\frac{1}{P(s)}\frac{L(s)}{1+L(s)} \tag{2.1.4}
\]

\[
T_1(s)\stackrel{\text{def}}{=}\frac{y(s)}{r_1(s)}=\frac{L(s)}{1+L(s)}\equiv T_{uf}(s) \tag{2.1.5}
\]

\[
T_{ud}(s)\stackrel{\text{def}}{=}\frac{u(s)}{d(s)}=\frac{-G(s)}{1+L(s)} \tag{2.1.6}
\]

\[
T_{yn}(s)=\frac{y(s)}{n(s)}=-\frac{L(s)}{1+L(s)} \tag{2.1.7}
\]

\subsection{Significado físico de las funciones de transferencia}
\begin{itemize}
  \item \textbf{\(T(s)\)} es la función de seguimiento entrada-salida en una estructura \textbf{TDOF}.
  \item \textbf{\(T_1(s)\)} es la función de seguimiento en una estructura \textbf{ODOF} (un solo grado de libertad), es decir, cuando no existe o no puede implementarse el prefiltro \(F(s)\).
  \item \textbf{\(T_d(s)\)} describe cómo una perturbación aplicada a la salida de la planta afecta a la salida del sistema.
  \item \textbf{\(T_{d1}(s)\)} hace lo mismo para una perturbación aplicada a la entrada de la planta.
  \item \textbf{\(S(s)\)} es la \textbf{función de sensibilidad} y es crucial porque también mide la sensibilidad del sistema frente a incertidumbre de la planta.
  \item \textbf{\(T_{un}(s)\)} cuantifica la amplificación del ruido del sensor hacia la entrada de la planta.
  \item \textbf{\(T_{yn}(s)\)} describe cuánto del ruido del sensor llega a la salida del sistema.
\end{itemize}

Una identidad fundamental es:

\[
T_1(s)+S(s)=1 \tag{2.1.8}
\]

Esta ecuación significa que, en un sistema \textbf{ODOF}, las dos funciones más importantes —la complementaria \(T_1(s)\) y la sensibilidad \(S(s)\)— \textbf{no pueden imponerse de manera independiente}.

También es útil escribir la salida y la entrada de la planta en términos de todas las señales externas:

\[
y(s)=T(s)r(s)+S(s)d(s)+P(s)S(s)d_1(s)-T_1(s)n(s) \tag{2.1.9}
\]

$$
u(s)=F(s)G(s)S(s)r(s)-G(s)S(s)d(s)+S(s)d_1(s)-G(s)S(s)n(s)
$$

De aquí se obtienen varias conclusiones de diseño:

\begin{enumerate}
  \item Para atenuar perturbaciones \(d\) y \(d_1\), se necesita que \(|S(j\omega)|\) sea pequeño en los rangos de frecuencia donde estas perturbaciones son significativas.
  \item Para evitar que el ruido del sensor domine la salida, se requiere que \(|T_1(j\omega)|\) sea pequeño en la banda donde el ruido está concentrado.
  \item Para mantener bajo el esfuerzo de control y evitar saturación, \(|G(j\omega)|\) debe ser pequeño, sobre todo a altas frecuencias.
\end{enumerate}

La consecuencia inmediata es que el diseño de \(L(s)=G(s)P(s)\) exige \textbf{compromisos}: conviene una gran magnitud de \(|L(j\omega)|\) en frecuencias bajas e intermedias, pero una magnitud pequeña en frecuencias altas, donde domina el ruido.



\section{2.2 Definición matemática de la retroalimentación}

\subsection{2.2.1 El origen de la teoría de la retroalimentación}
La teoría de la retroalimentación surge históricamente en el diseño de \textbf{amplificadores con realimentación}, con trabajos de \textbf{Nyquist, Black y Bode}. El objetivo original era reducir la sensibilidad de la función de transferencia frente a incertidumbres de hardware.

Para ilustrarlo, el autor considera \(G(s)=k\), y retoma la definición de sensibilidad de Bode:

$$
S_k^T(s)=\frac{\partial \ln T}{\partial \ln k}
=\frac{\partial T/T}{\partial k/k}
$$

Aplicada a la estructura de la Figura 2.1.1, se obtiene:

$$
S_k^T(s)=\frac{1}{1+kP(s)}=\frac{1}{1+L(s)}
$$

es decir,

$$
S_k^T(s)=S(s)
$$

La interpretación es muy importante: la incertidumbre relativa de la transferencia cerrada \(T(s)\) respecto a la incertidumbre relativa del elemento directo queda reducida por el factor \(1+L(s)\). En otras palabras, la retroalimentación reduce sensibilidad a variaciones paramétricas.

\subsection{2.2.2 Razón de retorno, diferencia de retorno e invariancia de su numerador}
Un sistema real puede contener varios lazos locales de retroalimentación. Sin embargo, al final puede llevarse a una forma canónica equivalente. Para analizar estructuras de múltiples lazos, el autor usa \textbf{grafos de flujo de señal} y la idea de \textbf{razón de retorno}.

La idea operativa es la siguiente:

\begin{itemize}
  \item se anulan las entradas externas;
  \item se “abre” el punto del sistema que se quiere usar como referencia;
  \item se inyecta una señal de prueba en el camino directo;
  \item se calcula la señal que regresa por el camino de retorno.
\end{itemize}

Con ello se define la \textbf{razón de retorno} y, a partir de ella, la \textbf{diferencia de retorno}. El resultado esencial es que, aunque la forma concreta de la diferencia de retorno dependa del elemento elegido como referencia, \textbf{el numerador de esa diferencia es siempre el mismo}: coincide con el \textbf{polinomio característico del sistema}.

En notación compacta, si \(\Delta\) es el determinante del grafo y \(\Delta^k\) el determinante asociado al elemento de referencia \(k\), entonces:

\[
F_k=\frac{\Delta}{\Delta^k} \tag{2.2.6}
\]

y el numerador de \(\Delta\) no depende del elemento escogido. Esa invariancia explica por qué los polos de lazo cerrado del sistema son independientes del punto desde el que se “mire” el lazo.

\insertimage[\label{img:grafo-flujo-senal}]{2024_06_29_804c621bf2d98037ffbeg-023_804_1124_194_772}{width=0.45\textwidth}{Representación mediante grafo de flujo de señal de una estructura con múltiples lazos.}

El ejemplo del libro con los elementos \(k_1\) y \(k_2\) muestra que las diferencias de retorno calculadas respecto a ambos elementos poseen \textbf{denominadores distintos}, pero \textbf{el mismo numerador}:

$$
D(s)=Js^2+k_2s+k_1
$$

\insertimage[\label{img:diferencia-retorno-k1-k2}]{2024_06_29_804c621bf2d98037ffbeg-023_330_1102_1612_2217}{width=0.45\textwidth}{Cálculo de la diferencia de retorno tomando como referencia \(k_1\) y \(k_2\).}

En la estructura canónica de la Figura 2.1.1, si el lazo se abre en la entrada de \(G(s)\), la razón de retorno es \(-L(s)\) y la diferencia de retorno es:

$$
F(s)=1+L(s)
$$

Este mismo término aparece en el denominador de todas las funciones de transferencia del sistema.

\subsection{2.2.3 Estabilidad asintótica e interna}

\subsubsection{Estabilidad asintótica}
En sistemas LTI racionales, una función de transferencia es \textbf{asintóticamente estable} si y solo si \textbf{todos sus polos están en el semiplano izquierdo abierto} del plano \(s\). Polos sobre el eje imaginario se consideran inestables en este contexto.

Eso implica dos cosas:

\begin{itemize}
  \item toda entrada acotada produce una salida acotada;
  \item la parte transitoria decae exponencialmente con el tiempo.
\end{itemize}

\subsubsection{Estabilidad interna}
La \textbf{estabilidad interna} es más exigente. El sistema es internamente estable si, para toda señal externa acotada \((r, d_1, d, n)\), también son acotadas todas las variables internas relevantes \((e, u)\) y la salida \(y\).

Equivale a exigir que \textbf{todas} las funciones de transferencia desde entradas externas hacia variables internas y de salida sean estables.

El resultado central del capítulo es:

\textbf{Teorema 2.2.} El sistema de retroalimentación es internamente estable si y solo si:

a. \(1+L(s)\) no tiene ceros en el semiplano derecho;\\
b. al formar \(L(s)=P(s)G(s)\), \textbf{no hay cancelaciones polo-cero en el semiplano derecho cerrado}.

Esto es crucial. No basta con que el numerador de \(1+L(s)\) sea estable. Puede ocurrir que la transferencia de mando \(y/r\) resulte estable, pero alguna transferencia interna como \(y/d_1\) o \(u/n\) siga siendo inestable debido a una cancelación inadmisible de un polo o cero inestable de la planta. Por eso, los polos y ceros inestables de la planta \textbf{deben permanecer explícitamente en \(L(s)\)} y no ser “cancelados” por el compensador.



\section{2.3 Definición de algunos problemas de control a partir de la ecuación básica de retroalimentación}
Todas las funciones de transferencia de interés comparten el mismo término \(1+L(s)\). En consecuencia, una vez que se diseña \(G(s)\), queda fijado el comportamiento de \(S(s)\), \(T_1(s)\), \(T_d(s)\), \(T_{d1}(s)\), \(T_{un}(s)\), etc. Solo \(T(s)\) conserva un grado adicional de manipulación mediante el prefiltro \(F(s)\).

El autor enfatiza que, en control clásico, el diseño suele realizarse en el \textbf{dominio de la frecuencia}. Allí se usan principalmente:

\begin{itemize}
  \item el criterio de \textbf{Nyquist};
  \item los diagramas de \textbf{Bode};
  \item la \textbf{carta de Nichols}.
\end{itemize}

\subsection{2.3.1 Consideraciones de estabilidad asintótica y relativa en el plano \(s\) y en el dominio \(\omega\)}

\subsubsection{El criterio de estabilidad de Nyquist}
Sea \(\Gamma_N\) el contorno de Nyquist que encierra todo el semiplano derecho, y sea \(\Gamma_L\) su imagen por la función \(L(s)\). El criterio de Nyquist establece que si la imagen de ese contorno rodea el punto crítico \([-1+j0]\) un número \(N\) de veces, y \(L(s)\) tiene \(p\) polos en el semiplano derecho, entonces el número de ceros \(z\) de \(1+L(s)\) en el semiplano derecho es:

\[
z=N+p \tag{2.3.1}
\]

La estabilidad de lazo cerrado se logra si y solo si \(z=0\), es decir, si:

$$
N=-p
$$

Dicho de otro modo: si \(L(s)\) tiene \(p\) polos inestables, el contorno de Nyquist debe rodear el punto crítico \([-1+j0]\) exactamente \(p\) veces en el sentido adecuado para compensarlos.

\insertimage[\label{img:criterio-nyquist-punto-critico}]{2024_06_29_804c621bf2d98037ffbeg-026_580_1125_94_808}{width=0.45\textwidth}{Interpretación gráfica del criterio de Nyquist y del punto crítico de estabilidad.}

\subsubsection{Estabilidad relativa: márgenes de ganancia y fase}
La \textbf{estabilidad relativa} mide qué tan cerca está el lazo abierto de la inestabilidad. Los indicadores clásicos son:

\begin{itemize}
  \item \textbf{Margen de fase (PM)}: cuánto retraso de fase adicional puede tolerarse en la frecuencia de cruce sin perder estabilidad.
  \item \textbf{Margen de ganancia (GM)}: cuánto puede aumentarse la ganancia en la frecuencia donde la fase vale \(-180^\circ\) antes de provocar inestabilidad.
\end{itemize}

El punto \([-1+j0]\) se denomina \textbf{punto crítico de estabilidad}.

\subsection{2.3.2 Diseño en el dominio de la frecuencia con técnicas de Bode}
Antes del diseño con Bode, el autor recuerda dos nociones básicas:

\begin{itemize}
  \item la \textbf{magnitud logarítmica};
  \item la medición en \textbf{decibelios}:
\end{itemize}

$$
20\log_{10}|P(j\omega)|
$$

Un diagrama de Bode representa, en escala semilogarítmica:

\begin{itemize}
  \item la magnitud en dB contra \(\log \omega\);
  \item la fase contra \(\log \omega\).
\end{itemize}

La gran ventaja es que el comportamiento de una transferencia racional puede obtenerse sumando y restando aritméticamente las contribuciones de polos y ceros. Así, una vez conocido \(P(j\omega)\), el diseñador elige \(G(j\omega)\) para obtener un \(L(j\omega)\) que cumpla con:

\begin{itemize}
  \item coeficientes de error en estado estacionario;
  \item frecuencia de cruce \(\omega_{co}\);
  \item márgenes de fase y ganancia adecuados;
  \item estabilidad del sistema en lazo cerrado.
\end{itemize}

\insertimage[\label{img:bode-margenes-frecuencias}]{2024_06_29_804c621bf2d98037ffbeg-027_960_1000_136_817}{width=0.45\textwidth}{Definición de \(\omega_{co}\), \(\omega_{GM}\), margen de fase y margen de ganancia en el diagrama de Bode.}

\subsection{2.3.3 La carta de Nichols y sus características especiales}
La desventaja principal del diseño solo con Bode es que no muestra de forma explícita el comportamiento de la transferencia cerrada. La \textbf{carta de Nichols} sí lo hace.

La transferencia unitaria cerrada es:

\[
T_1(s)=\frac{L(s)}{1+L(s)} \tag{2.3.3}
\]

En la carta de Nichols se dibuja \(L(j\omega)\) usando:

\begin{itemize}
  \item magnitud de lazo abierto en dB sobre el eje vertical;
  \item fase de lazo abierto en grados sobre el eje horizontal.
\end{itemize}

Sobre ese plano se superponen curvas de nivel de magnitud constante de \(|T_1(j\omega)|\). Esto permite ver directamente:

\begin{itemize}
  \item el valor pico de \(|T_1|\);
  \item los márgenes de estabilidad relativa;
  \item el compromiso entre ancho de banda y sobreoscilación.
\end{itemize}

\insertimage[\label{img:carta-nichols}]{2024_06_29_804c621bf2d98037ffbeg-027_910_1148_996_2176}{width=0.45\textwidth}{Carta de Nichols con contornos de ganancia constante de \(|L/(1+L)|\).}

\subsubsection{Ubicación del punto crítico y lectura de márgenes}
En la carta de Nichols, el punto crítico de Nyquist aparece en:

$$
[0\ \text{dB},\ -180^\circ]
$$

y desde ahí se identifican de forma inmediata los márgenes de fase y de ganancia.

\subsubsection{Carta de Nichols invertida}
Como

$$
S(j\omega)=\frac{1}{1+L(j\omega)}
$$

y si se define \(l(s)=1/L(s)\), entonces:

\[
S(j\omega)=\frac{l(j\omega)}{1+l(j\omega)} \tag{2.3.4}
\]

Esto lleva a la \textbf{carta de Nichols invertida}, cuyas curvas corresponden a niveles constantes de \(|S(j\omega)|\). Geométricamente, esas curvas se obtienen al rotar las curvas de \(|T_1|\) alrededor del punto crítico. Esta herramienta es muy útil para estudiar rechazo de perturbaciones y sensibilidad.

\insertimage[\label{img:carta-nichols-invertida}]{2024_06_29_804c621bf2d98037ffbeg-028_910_1147_170_2212}{width=0.45\textwidth}{Carta de Nichols invertida, con contornos de \(|S|=|1/(1+L)|\), superpuesta sobre la carta de Nichols estándar.}

\subsubsection{¿Por qué Bode y Nichols son superiores al diseño directo en el plano complejo?}
El autor remarca dos razones prácticas:

\begin{enumerate}
  \item En Bode y Nichols, magnitudes y fases se \textbf{suman} aritméticamente; en el plano complejo, el moldeado de \(L(j\omega)\) es mucho menos intuitivo.
  \item A altas frecuencias, donde \(|L(j\omega)|\) es muy pequeño, la escala logarítmica de Bode/Nichols permite trabajar con más claridad y precisión.
\end{enumerate}

\subsection{2.3.4 Distinción entre sistemas ODOF y TDOF}
La diferencia esencial es la existencia o no del prefiltro \(F(s)\).

\begin{itemize}
  \item En un sistema \textbf{TDOF}, \(G(s)\) puede usarse para resolver robustez, rechazo de perturbaciones y márgenes de estabilidad, mientras que \(F(s)\) se utiliza para ajustar la transferencia entrada-salida \(T(s)\) según la especificación de seguimiento.
  \item En un sistema \textbf{ODOF}, solo existe \(G(s)\). En ese caso, el diseño debe comprometer simultáneamente propiedades de \(L(s)\) y de \(T(s)\), porque no se pueden especificar por separado.
\end{itemize}

El autor pone dos ejemplos conceptuales:

\begin{itemize}
  \item un sistema de control de vuelo de una aeronave, donde \(F(s)\) sí puede implementarse físicamente;
  \item un sistema de seguimiento por radar, donde solo el error es medible y por tanto no se puede introducir un prefiltro físico independiente.
\end{itemize}

El capítulo subraya una consecuencia importante: en un sistema TDOF se puede usar \(F(s)\) para cancelar dinámicas subamortiguadas de \(T_1(s)\) y obtener la respuesta temporal deseada sin alterar directamente la robustez del lazo principal. En un sistema ODOF esa separación no es posible.

\subsection{2.3.5 Diseño en frecuencia de sistemas ODOF}
El diseño comienza moldeando \(L(j\omega)\). Los parámetros más usados son:

\begin{itemize}
  \item \textbf{frecuencia de cruce} \(\omega_{co}\), donde \(|L(j\omega)|=1\);
  \item \textbf{frecuencia de margen de ganancia} \(\omega_{GM}\), donde \(\arg L(j\omega)=-180^\circ\);
  \item ancho de banda de lazo cerrado;
  \item rechazo de perturbaciones;
  \item sensibilidad al ruido.
\end{itemize}

\subsubsection{Frecuencia de cruce y estabilidad condicional}
La frecuencia de cruce está relacionada con la rapidez del sistema, aunque no es igual al ancho de banda de lazo cerrado. Si el lazo tiene comportamientos más complejos —por ejemplo, varios cruces o polos inestables—, los diagramas de Bode por sí solos no bastan, y la estabilidad debe revisarse con Nyquist.

\subsubsection{Desempeño de lazo cerrado y ancho de banda \(\omega_{-3dB}\)}
El ancho de banda de \(|T_1(j\omega)|\) surge como resultado del diseño del lazo abierto; en un sistema ODOF no suele especificarse de manera independiente. Está fuertemente relacionado con \(\omega_{co}\), pero no coincide exactamente con él.

\subsubsection{Rechazo de perturbaciones}
Para obtener buen rechazo, es deseable que \(|L(j\omega)|\) sea grande en el rango donde se concentra la perturbación. Eso hace pequeño a \(|S(j\omega)|\).

\subsubsection{Loopshaping}
El moldeado de \(L(j\omega)\) se realiza con redes:

\begin{itemize}
  \item \textbf{lead-lag},
  \item \textbf{lag-lead},
  \item combinaciones más complejas,
  \item y polos lejanos (“far-off poles”) para reducir ruido a altas frecuencias.
\end{itemize}

Los ejemplos del libro ilustran cómo usar estas redes para cumplir simultáneamente con especificaciones como:

\begin{itemize}
  \item coeficiente de error de velocidad \(K_v\),
  \item margen de fase,
  \item frecuencia de cruce deseada,
  \item y limitación del ruido.
\end{itemize}

Los mensajes principales de los ejemplos son:

\begin{itemize}
  \item una red \textbf{lead-lag} puede aumentar fase cerca de \(\omega_{co}\) y mejorar los márgenes;
  \item una red \textbf{lag-lead} puede ajustar la magnitud sin sacrificar demasiado la fase;
  \item un \textbf{polo lejano adicional} ayuda a atenuar la amplificación del ruido del sensor.
\end{itemize}

\begin{images}[\label{img:loopshaping-bode-nichols}]{Ejemplos de loopshaping en Bode y Nichols mediante redes lead-lag y lag-lead.}
  \addimage{2024_06_29_804c621bf2d98037ffbeg-032_384_1040_183_2248}{width=0.45\textwidth}{2.3.8}
  \imagesnewline
  \addimage{2024_06_29_804c621bf2d98037ffbeg-032_362_1014_605_2243}{width=0.45\textwidth}{2.3.9}
  \imagesnewline
  \addimage{2024_06_29_804c621bf2d98037ffbeg-032_835_1026_1065_2237}{width=0.45\textwidth}{2.3.10}
  \imagesnewline
  \addimage{2024_06_29_804c621bf2d98037ffbeg-033_784_1024_198_2237}{width=0.45\textwidth}{2.3.11a}
  \imagesnewline
  \addimage{2024_06_29_804c621bf2d98037ffbeg-033_812_1036_1103_2214}{width=0.45\textwidth}{2.3.11b}
\end{images}

\subsubsection{Limitaciones de los sistemas ODOF}
En los sistemas de un solo grado de libertad, la transferencia de seguimiento y las propiedades del lazo no pueden fijarse en forma independiente. El ancho de banda de \(|T_1|\) y su forma son consecuencia indirecta del diseño del lazo.

\subsection{2.3.6 Diseño en frecuencia de sistemas TDOF}
Si el prefiltro \(F(s)\) puede implementarse, entonces la transferencia total:

$$
T(s)=T_1(s)F(s)
$$

puede moldearse casi a voluntad una vez que \(T_1(s)\) ya fue obtenido mediante el diseño de \(L(s)\). Eso permite satisfacer exactamente una especificación adicional de seguimiento sin arruinar las propiedades del lazo principal.

En el ejemplo 2.3.4, el libro muestra cómo, partiendo de un lazo ya diseñado, se agrega un prefiltro para obtener un ancho de banda cerrado deseado de aproximadamente:

$$
\omega_{-3dB}\approx 1\ \text{rad/s}
$$

La idea conceptual es clara: en TDOF, el diseño del controlador \(G(s)\) y el del prefiltro \(F(s)\) son procesos \textbf{separables}.



\section{2.4 Compromisos de diseño en sistemas de retroalimentación}
La necesidad de compromisos aparece porque el diseñador busca simultáneamente:

\begin{itemize}
  \item buen seguimiento,
  \item fuerte rechazo de perturbaciones,
  \item baja sensibilidad a incertidumbre,
  \item poco ruido amplificado,
  \item y esfuerzo de control limitado.
\end{itemize}

Todo eso no puede maximizarse al mismo tiempo.

\subsection{2.4.1 El problema de amplificación del ruido del sensor}
La transferencia del ruido del sensor a la entrada de la planta es:

$$
|T_{un}(j\omega)|=\left|\frac{G(j\omega)}{1+L(j\omega)}\right|
=\left|\frac{1}{P(j\omega)}\frac{L(j\omega)}{1+L(j\omega)}\right|
$$

El autor divide el espectro de frecuencias en tres zonas.

\subsubsection{a) Frecuencias bajas: \(|L(j\omega)|\gg 1\)}
\[
|T_{un}(j\omega)|\approx \left|\frac{1}{P(j\omega)}\right| \tag{2.4.2}
\]

En esta región, la amplificación del ruido está determinada principalmente por la planta.

\subsubsection{b) Frecuencias altas: \(|L(j\omega)|\ll 1\)}
\[
|T_{un}(j\omega)|\approx |G(j\omega)|
=\left|\frac{L(j\omega)}{P(j\omega)}\right| \tag{2.4.3}
\]

Esta es la región más importante para el ruido, porque allí suele concentrarse el espectro del sensor.

\subsubsection{c) Frecuencias intermedias: \(|L(j\omega)|\approx 1\)}
En esta zona el valor de \(|T_{un}|\) depende de qué tan cerca pase \(L(j\omega)\) del punto crítico. Si el diseño evita picos excesivos de \(|T_1|\) y \(|S|\), entonces también se limita la amplificación de ruido en esta banda.

\insertimage[\label{img:ruido-sensor-frecuencia}]{2024_06_29_804c621bf2d98037ffbeg-035_942_1174_172_2181}{width=0.45\textwidth}{Generación de amplificación de ruido del sensor en el dominio de la frecuencia.}

La conclusión importante es esta: para que el valor RMS del ruido amplificado sea finito, es necesario que \(|T_{un}(j\omega)|\to 0\) cuando \(\omega\to\infty\). Por tanto:

\begin{itemize}
  \item \(G(s)\) debe ser \textbf{estrictamente propia};
  \item la pendiente de \(|L(j\omega)|\) a altas frecuencias debe ser más negativa que la de \(|P(j\omega)|\).
\end{itemize}

\subsection{2.4.2 Cálculo RMS de \(|T_{un}(j\omega)|\)}
El libro muestra, mediante un ejemplo, cómo estimar gráficamente el valor RMS del ruido amplificado integrando el área bajo \(|T_{un}(j\omega)|^2\). En el ejemplo presentado:

\begin{itemize}
  \item cálculo gráfico: \(\approx 201.6\),
  \item cálculo exacto en Matlab: \(\approx 203.74\).
\end{itemize}

\begin{images}[\label{img:calculo-rms-ruido}]{Cálculo gráfico de la amplificación RMS del ruido del sensor.}
  \addimage{2024_06_29_804c621bf2d98037ffbeg-036_916_1147_176_748}{width=0.45\textwidth}{2.4.2}
  \imagesnewline
  \addimage{2024_06_29_804c621bf2d98037ffbeg-036_901_1151_188_2179}{width=0.45\textwidth}{2.4.3}
\end{images}

El punto didáctico del ejemplo es que una magnitud de \(L(j\omega)\) demasiado alta en una banda extensa puede dar lugar a una amplificación RMS muy severa.

\subsection{2.4.3 Optimización de la función de transmisión de lazo \(L(j\omega)\)}
En el rango donde \(|L(j\omega)|>1\), aumentar \(|L|\) mejora en general los beneficios de la retroalimentación: baja sensibilidad, mejor rechazo de perturbaciones y mejor seguimiento.

Sin embargo, no puede aumentarse indiscriminadamente. El autor muestra que el criterio práctico consiste en \textbf{explotar al máximo el retraso de fase permitido} sin arruinar el margen de fase y sin acercarse demasiado al punto crítico.

La idea es usar una red apropiada para aumentar la magnitud de \(L(j\omega)\) debajo de \(\omega_{co}\), manteniendo todavía un margen de fase satisfactorio.

\begin{images}[\label{img:optimizacion-lazo-laglead}]{Implementación del criterio de optimización del lazo y uso conceptual de una red lag-lead.}
  \addimage{2024_06_29_804c621bf2d98037ffbeg-037_924_1142_176_746}{width=0.45\textwidth}{2.4.4}
  \imagesnewline
  \addimage{2024_06_29_804c621bf2d98037ffbeg-037_676_664_1262_986}{width=0.45\textwidth}{2.4.5}
\end{images}



\section{2.5 Loopshaping basado en optimización con norma \(H_\infty\)}
Hasta aquí el diseño se presentó con herramientas clásicas moldeando directamente \(L(s)\). Ahora el autor reescribe las funciones cerradas en términos de \(S(s)\) y \(T_1(s)\):

\[
T(s)=F(s)T_1(s) \tag{2.5.1}
\]

\[
T_{d1}(s)=P(s)S(s) \tag{2.5.2}
\]

\[
T_d(s)=S(s) \tag{2.5.3}
\]

\[
T_{un}(s)=-G(s)S(s)=T_{ud}(s) \tag{2.5.4}
\]

$$
T_1(s)=1-S(s)=G(s)P(s)S(s)=-T_{yn}(s)
$$

Además, en frecuencias bajas y altas se cumplen aproximaciones útiles:

$$
|L(j\omega)|\gg 1 \Rightarrow S(j\omega)\approx \frac{1}{L(j\omega)},\quad T_1(j\omega)\approx 1
$$

$$
|L(j\omega)|\ll 1 \Rightarrow S(j\omega)\approx 1,\quad T_1(j\omega)\approx |L(j\omega)|
$$

En la región intermedia, donde \(|L|\approx 1\), ya no valen simplificaciones tan directas, y es justamente ahí donde pueden aparecer picos indeseables de \(|S|\) y \(|T_1|\).

\subsection{2.5.1 Definición de las normas \(H_\infty\) y \(H_2\)}
La norma \(H_\infty\) de una función estable y propia \(F(s)\) se define como:

\[
\|F(s)\|_\infty \stackrel{\text{def}}{=} \max_\omega |F(j\omega)| \tag{2.5.6}
\]

o, si el máximo no existe de forma estricta,

\[
\|F(s)\|_\infty = \sup_\omega |F(j\omega)| \tag{2.5.6a}
\]

Es, en esencia, el \textbf{valor pico} del diagrama de magnitud de Bode.

La norma \(H_2\) se define como:

$$
\|F(j\omega)\|_2=
\left[
\frac{1}{2\pi}\int_{-\infty}^{\infty}|F(j\omega)|^2\,d\omega
\right]^{1/2}
$$

También vale la propiedad:

$$
\|F_1(s)F_2(s)\|_\infty \le \|F_1(s)\|_\infty \,\|F_2(s)\|_\infty
$$

El capítulo recuerda además:

\begin{itemize}
  \item \(\|F\|_2\) es finita si \(F\) es estrictamente propia y no tiene polos sobre el eje imaginario;
  \item \(\|F\|_\infty\) es finita si \(F\) es propia y no tiene polos sobre el eje imaginario.
\end{itemize}

\subsubsection{Interpretación física}
La norma \(H_2\) aparece cuando se quiere minimizar energía promedio o esperada de señales como el error y el esfuerzo de control. En cambio, la norma \(H_\infty\) aparece de manera natural cuando interesa controlar el \textbf{peor caso}.

Por ejemplo, si se quiere penalizar simultáneamente error y esfuerzo de control, puede definirse un índice del tipo:

$$
PI = E\left\{\|e\|_2^2+\beta^2\|u\|_2^2\right\}
$$

que conduce a una formulación \(H_2\). Si lo que interesa es limitar el peor caso sobre perturbaciones con energía acotada, entonces aparece la formulación \(H_\infty\).

\subsection{2.5.2 Requisitos básicos de estabilidad relativa}
El autor relaciona los márgenes clásicos con los picos máximos de \(|T_1|\) y \(|S|\). La idea es muy útil porque traduce márgenes geométricos a especificaciones directas sobre magnitudes cerradas.

\begin{itemize}
  \item Si \(|T_1(j\omega)|_{\max}\le 3\ \text{dB}\), queda garantizado aproximadamente que:

  \begin{itemize}
    \item \(PM>45^\circ\),
    \item \(GM>4\ \text{dB}\).
  \end{itemize}
  \item Si \(|S(j\omega)|_{\max}\le 3\ \text{dB}\), queda garantizado aproximadamente que:

  \begin{itemize}
    \item \(PM>45^\circ\),
    \item \(GM>11\ \text{dB}\).
  \end{itemize}
\end{itemize}

En particular, para \(|S(j\omega)|_{\max}\), el texto deduce:

\[
GM \ge \frac{|S(j\omega)|_{\max}}{|S(j\omega)|_{\max}-1} \tag{2.5.12}
\]

y

\[
PM \ge 2\sin^{-1}\left[\frac{1}{2|S(j\omega)|_{\max}}\right] \tag{2.5.13}
\]

Esto permite traducir restricciones de robustez y estabilidad relativa a cotas directamente sobre el pico de la sensibilidad.

\subsection{2.5.3 Sensibilidad ponderada}
Como \(S(s)\) interviene en:

\begin{itemize}
  \item sensibilidad a incertidumbre,
  \item error de seguimiento,
  \item rechazo de perturbaciones,
  \item estabilidad relativa,
\end{itemize}

se desea que su magnitud sea pequeña en la banda adecuada. Pero esa exigencia no debe ser uniforme en todas las frecuencias; debe depender del problema.

Por eso se introduce una función de ponderación \(W_S(s)\) tal que:

\[
|S(j\omega)|<\frac{1}{|W_S(j\omega)|},\qquad \forall \omega \tag{2.5.15}
\]

equivalentemente,

\[
\|W_S(j\omega)S(j\omega)\|_\infty<1 \tag{2.5.16}
\]

Esto quiere decir que la magnitud de \(S\) debe quedar por debajo de una “envolvente” definida por \(1/|W_S|\).

Además, como \(S=1/(1+L)\), de ahí se obtiene:

\[
|1+L(j\omega)|>|W_S(j\omega)| \tag{2.5.17}
\]

lo que puede interpretarse geométricamente como una restricción sobre la distancia de \(L(j\omega)\) al punto crítico \(-1\).

\insertimage[\label{img:especificaciones-s-circulos}]{2024_06_29_804c621bf2d98037ffbeg-041_839_911_1134_848}{width=0.45\textwidth}{Interpretación geométrica de las especificaciones sobre \(S(j\omega)\) en términos de círculos alrededor del punto crítico.}

\subsection{2.5.4 Sensibilidad mixta}
En la práctica no basta con imponer una cota a \(S\). También interesa limitar, por ejemplo:

\begin{itemize}
  \item la transferencia de ruido a la entrada de la planta \(T_{un}\),
  \item la transferencia complementaria \(T_1\).
\end{itemize}

Así surgen especificaciones simultáneas del tipo:

\[
\|W_S S\|_\infty<1,\qquad
\|W_{un}T_{un}\|_\infty<1,\qquad
\|W_1T_1\|_\infty<1 \tag{2.5.19}
\]

Como en un sistema ODOF solo existe un controlador \(G(s)\), tales especificaciones no pueden cumplirse todas de forma exacta e independiente. De ahí la idea de \textbf{sensibilidad mixta}, reuniéndolas en un solo vector de desempeño:

\[
\mathbf{PV}=
\begin{bmatrix}
W_S S\\
W_{un}T_{un}\\
W_1T_1
\end{bmatrix}
\tag{2.5.20}
\]

y minimizando:

\section{\$\$\\
|\textbackslash mathbf\{PV\}(j\textbackslash omega)|\_\textbackslash infty}
\textbackslash max\_\textbackslash omega\\
\textbackslash sqrt\{|W\_SS|\textsuperscript{2+|W\_\{un\}T\_\{un\}|}2+|W\_1T\_1|\^{}2\}\\
<1 \textbackslash tag\{2.5.21\}

\$\$

El problema de diseño se formula entonces como:

\[
\underset{G(s)}{\text{minimizar}}\ \|\mathbf{PV}(G)\|_\infty \tag{2.5.22}
\]

El autor advierte que esta formulación es conservadora: con \(n\) requisitos mezclados, el error máximo respecto al cumplimiento individual puede llegar a un factor \(\sqrt{n}\).

\subsection{2.5.5 El problema regulador estándar \(H_\infty\)}
La sensibilidad mixta es un caso particular del \textbf{problema regulador estándar \(H_\infty\)}. Para formularlo, se construye una \textbf{planta generalizada} \(M(s)\) que incorpora la planta física, las ponderaciones y las señales exógenas.

En la estructura correspondiente aparecen tres salidas ponderadas, por ejemplo:

$$
z_1=-S W_S\,d
$$

$$
z_2=-G S W_{un}\,n=-T_{un}W_{un}\,n
$$

\[
z_3=T_1 W_1\,r_1 \tag{2.5.23}
\]

y se definen:

$$
\mathbf{w}(s)=\text{col}[d(s),n(s),r_1(s)]
$$

$$
\mathbf{z}(s)=\text{col}[z_1(s),z_2(s),z_3(s)]
$$

donde \(\mathbf{z}\) es el vector de errores ponderados que se desea minimizar.

\begin{images}[\label{img:sensibilidad-mixta-regulador}]{Estructura de sensibilidad mixta y forma canónica del problema regulador estándar \(H_\infty\).}
  \addimage{2024_06_29_804c621bf2d98037ffbeg-042_496_1152_1462_741}{width=0.45\textwidth}{2.5.2}
  \imagesnewline
  \addimage{2024_06_29_804c621bf2d98037ffbeg-042_358_469_1580_2489}{width=0.45\textwidth}{2.5.3}
\end{images}

En forma compacta:

\[
\mathbf{z}=\mathbf{T}_{wz}\mathbf{w} \tag{2.5.28}
\]

y el problema general consiste en encontrar el controlador que minimice:

\[
\underset{\mathbf{G}(s)}{\text{minimizar}}\ \|\mathbf{T}_{wz}\|_\infty \tag{2.5.29}
\]

\subsection{2.5.6 Selección de funciones de ponderación}
La elección de \(W_S(s)\), \(W_1(s)\) y \(W_{un}(s)\) es el paso inicial y más delicado del diseño. Deben codificar en forma matemática los requisitos del cliente o del problema físico:

\begin{itemize}
  \item espectro y atenuación requerida para perturbaciones;
  \item error estacionario máximo permitido;
  \item ancho de banda y sobrepico admisibles de seguimiento;
  \item ruido del sensor o restricciones especiales sobre el esfuerzo de control.
\end{itemize}

El autor aclara que \(W_{un}(s)\) puede tomarse inicialmente igual a 1 si no existe una exigencia particular, y refinarse después.

\subsection{2.5.7 Solución \(H_\infty\) del problema de sensibilidad mixta}
El capítulo resuelve un ejemplo completo.

\subsubsection{Ejemplo 2.5.1}
La planta es:

$$
P(s)=\frac{5}{(s+1)(s+0.1)}
$$

y las especificaciones son:

\begin{enumerate}
  \item \(K_v=10\);
  \item ancho de banda de \(S\): \(\omega_{S,-3dB}=5\ \text{rad/s}\);
  \item \(PM>40^\circ\), equivalente a \(|S|_{\max}<1.46\approx 3.3\ \text{dB}\);
  \item \(|T_1|_{\max}<1.3\approx 2.3\ \text{dB}\);
  \item ancho de banda de \(|T_1|\): \(\omega_{-3dB}=8\ \text{rad/s}\);
  \item mantener lo más pequeña posible la amplificación del ruido del sensor.
\end{enumerate}

\subsubsection{Etapa 1: definición de \(W_S(s)\)}
A partir de \(K_v=10\) y del requisito de ancho de banda/sobrepico para \(S\), el texto propone:

$$
\frac{1}{W_S(s)}=
\frac{s(1+s/2.13)}{10(1+s/5.58)^2}
$$

\subsubsection{Etapa 2: definición de \(W_1(s)\)}
Para imponer el pico y el ancho de banda de \(T_1\), se elige:

$$
\frac{1}{W_1(s)}=
\frac{1.3}{(1+s/6)(1+s/20)(1+s/30)}
$$

\subsubsection{Etapa 3: definición de \(W_{un}(s)\)}
En este ejemplo se toma simplemente:

$$
W_{un}(s)=1
$$

\subsubsection{Etapa 4: solución con \(H_\infty\)}
Usando el algoritmo \textbf{hinfsyn} de \(\mu\)-Tools, el autor obtiene un compensador subóptimo de orden elevado:

$$
G(s)=
\frac{80000(s+100)\left(s^2+200s+100.16\right)(s+3)(s+1)(s+0.1)}
{(s+6091)\left(s^2+209s+10972\right)(s+91.17)(s+13.74)(s+2.13)(s+0.001)}
$$

\textbf{[Espacio para la Tabla 2.5.1]}\\
\textbf{Tabla 2.5.1.} Script de Matlab que construye la planta generalizada y aplica el algoritmo \texttt{hinfsyn}.

\subsubsection{Etapa 5: evaluación y comparación con diseño clásico}
El mismo problema se resuelve también por métodos clásicos, obteniéndose un compensador mucho más simple:

$$
G_{\text{clas}}(s)=\frac{2000(s+0.1)(s+1)}{s(s+20)(s+50)}
$$

El resultado comparativo es interesante:

\begin{itemize}
  \item ambas soluciones cumplen razonablemente bien;
  \item la solución clásica presenta anchos de banda algo mayores en \(S\) y \(T_1\);
  \item la solución \(H_\infty\) produce una amplificación RMS del ruido del sensor significativamente mayor:
\end{itemize}

$$
|T_{un}|_{H_\infty,\ RMS}\approx 767.431
$$

frente a

$$
|T_{un}|_{\text{clas},\ RMS}\approx 181.2
$$

La explicación es que el lazo abierto obtenido por \(H_\infty\) cae con una pendiente menos negativa en alta frecuencia. El autor recomienda, por ello, \textbf{afinar después la solución \(H_\infty\) con criterios clásicos}, agregando polos lejanos si es necesario para mejorar la atenuación del ruido del sensor sin degradar el desempeño en la banda útil.

\begin{images}[\label{img:comparacion-clasico-hinf}]{Comparación entre diseño clásico y diseño \(H_\infty\): sensibilidad, transferencia complementaria, lazo abierto y respuestas temporales.}
  \addimage{2024_06_29_804c621bf2d98037ffbeg-045_490_1159_170_742}{width=0.45\textwidth}{2.5.4}
  \imagesnewline
  \addimage{2024_06_29_804c621bf2d98037ffbeg-045_442_1139_1475_743}{width=0.45\textwidth}{2.5.5}
  \imagesnewline
  \addimage{2024_06_29_804c621bf2d98037ffbeg-045_848_1080_179_2237}{width=0.45\textwidth}{2.5.6}
  \imagesnewline
  \addimage{2024_06_29_804c621bf2d98037ffbeg-045_894_1108_1104_2208}{width=0.45\textwidth}{2.5.7}
  \imagesnewline
  \addimage{2024_06_29_804c621bf2d98037ffbeg-046_848_1121_179_770}{width=0.45\textwidth}{2.5.8}
\end{images}

Una observación final muy importante: si el prefiltro \(F(s)\) es físicamente implementable, la \(|T_1(j\omega)|\) obtenida con el diseño \(H_\infty\) también puede ajustarse después para satisfacer una especificación completa de seguimiento:

$$
|T(j\omega)|=|T_1(j\omega)|\,|F(j\omega)|
$$

\section{2.6 Loopshaping clásico versus loopshaping moderno con \(H_\infty\)}

% In the present chapter, feedback control was approached by classical and modern design techniques. In both design paradigms, a control network $G(s)$ was sought to shape the loop transmission function $L(s)=P(s) G(s)$, so that initially specified closed-loop performances are achieved. Some major differences exist between both approaches. In the classical approach, the closed-loop specifications are achieved by directly shaping the loop transmission function $L(j \omega)$ by Nyquist/Bode/Nichols oriented design procedures. With this approach, the design process is very transparent to the designer; by using the Bode plots and the Nichols chart, both open- and closedloop properties are under the designer's control during the loopshaping process. On the other hand, the modern $H_{\infty}$-feedback control design procedure consists in defin-ing weighting functions for the closed-loop performance specifications, and, by using $H_{\infty}$ optimization algorithms to obtain the control network $G(s)$. Thus, a large part of the design effort is transferred to the computer. The commercial $H_{\infty}$-optimization algorithms that are available change iteratively the parameters of the assumed $G(s)$ until a satisfactory solution is obtained. However, if the final outcome is not exactly what the designer expected to obtain, some of the weighting functions can be altered in a logical way so that an ameliorated solution will be achieved.

% When SISO feedback systems are designed, comparable solutions from system point of view are obtained with both approaches. However, for the MIMO case, the classical approach is less efficient and also less transparent to the designer, because shaping of any individual loop out of the $n$ loops to be designed, influences the previously obtained feedback properties of the remaining loops. This point will be clarified in Chapter 6, where the MIMO feedback design is treated.

En el presente capítulo se abordó el control por retroalimentación mediante técnicas de diseño clásicas y modernas. En ambos paradigmas de diseño, se buscó una red de control \(G(s)\) para moldear la función de transmisión de lazo \(L(s)=P(s)G(s)\), de modo que se cumplieran especificaciones de lazo cerrado inicialmente especificadas. Existen algunas diferencias importantes entre ambos enfoques. En el enfoque clásico, las especificaciones de lazo cerrado se logran moldeando directamente la función de transmisión de lazo \(L(j\omega)\) mediante procedimientos de diseño orientados a Nyquist/Bode/Nichols. Con este enfoque, el proceso de diseño es muy transparente para el diseñador; al usar los diagramas de Bode y el gráfico de Nichols, tanto las propiedades del lazo abierto como del lazo cerrado están bajo el control del diseñador durante el proceso de loopshaping. Por otro lado, el procedimiento moderno de diseño de control por retroalimentación \(H_\infty\) consiste en definir funciones de ponderación para las especificaciones de desempeño del lazo cerrado y, utilizando algoritmos de optimización \(H_\infty\), obtener la red de control \(G(s)\). Así, gran parte del esfuerzo de diseño se transfiere a la computadora. Los algoritmos comerciales de optimización \(H_\infty\) disponibles cambian iterativamente los parámetros del supuesto \(G(s)\) hasta que se obtiene una solución satisfactoria. Sin embargo, si el resultado final no es exactamente lo que el diseñador esperaba obtener, algunas de las funciones de ponderación pueden ser alteradas de manera lógica para lograr una solución mejorada.

Cuando se diseñan sistemas de retroalimentación SISO, se obtienen soluciones comparables desde el punto de vista del sistema con ambos enfoques. Sin embargo, para el caso MIMO, el enfoque clásico es menos eficiente y también menos transparente para el diseñador, porque moldear cualquiera de los lazos individuales de los \(n\) lazos a diseñar influye en las propiedades de retroalimentación previamente obtenidas de los lazos restantes. Este punto se aclarará en el Capítulo 6, donde se trata el diseño de retroalimentación MIMO.

\begin{sourcecode}{matlab}{Programa Matlab para diseñar el controlador subóptimo \(H_\infty\) \(G(s)\).}
% Hinf solution of Example 2.5.1; mu251.m; 20 december 1997
% Plant G = 5 / (s+0.1)(s+1)
% Ws = 10(1+s/5.58)^2 / ((s+0.001)(1+s/2.13))
% Wt = (1+s/6)(1+s/20)(1+s/30) / (1.3(1+s/100)^3)
% Wu = 1
% computation of G, Wt and Ws
ng = [1];
dg = [1 1.1 0.1];
G = nd2sys(ng, dg, 5); % plant
nws = 10 * [1/5.58^2 2/5.18 1];
dws = [1/2.13 1+0.001/2.13 0.001];
Ws = nd2sys(nws, dws, 1.0);
nwt = 213.7 * [1569003600];
dwt = [13003 c+4 le +6];
Wt = nd2sys(nwt, dwt, 1/1.3);
Wu = 1;
% end
% preparation of generalized plant P
systemnames = 'G Wt Wu Ws';
inputvar = '[w(1); u(1)]';
outputvar = '[Wt; Wu; Ws; w - G]';
input_to_G = '[u]';
input_to_Wt = '[G]';
input_to_Wu = '[u]';
input_to_Ws = '[G-u]';
sysoutname = 'P';
cleanupsysic = 'yes';
sysic;
% end
% start of Hinf algorithm
nmeas = 1;
nu = 1;
gmn = 0.5;
gmx = 20;
tol = 0.001;
[khinf, ghinf, gopt] = hinfsyn(P, nmeas, nu, gmn, gmx, tol);
% end
% calculation of suboptimal compensator
[a, b, c, d] = unpck(khinf);
[numg, deng] = ss2tf(a, b, c, d);
roots(numg)
roots(deng)
end
\end{sourcecode}

\section{2.7 Resumen}
Este capítulo establece las bases del diseño robusto en frecuencia para sistemas SISO.

\begin{enumerate}
  \item Se definen las transferencias fundamentales del problema de retroalimentación: \(T\), \(T_1\), \(S\), \(T_d\), \(T_{d1}\), \(T_{un}\), etc.
  \item Se explica la noción de sensibilidad de Bode y la importancia central del término \(1+L(s)\).
  \item Se distingue entre \textbf{estabilidad asintótica} y \textbf{estabilidad interna}, subrayando que no se permiten cancelaciones de polos o ceros en el semiplano derecho.
  \item Se introducen las herramientas clásicas de diseño: \textbf{Nyquist, Bode, Nichols y Nichols invertida}.
  \item Se aclara la diferencia entre sistemas \textbf{ODOF} y \textbf{TDOF}, y el papel del prefiltro \(F(s)\).
  \item Se muestra que el diseño del lazo está lleno de \textbf{compromisos} entre robustez, seguimiento, rechazo de perturbaciones, saturación y amplificación de ruido.
  \item Finalmente, se formula el mismo problema en lenguaje moderno mediante \textbf{sensibilidad ponderada}, \textbf{sensibilidad mixta} y el \textbf{regulador estándar \(H_\infty\)}, cerrando con una comparación práctica entre una solución clásica y una obtenida por optimización.
\end{enumerate}

La idea más importante del capítulo es que el diseño de control por retroalimentación no consiste en “hacer grande el lazo” sin más, sino en darle a \(L(j\omega)\) una forma cuidadosamente balanceada en frecuencia, de modo que el sistema sea \textbf{estable, robusto, realizable y físicamente sensato}.


\input{example} % Ejemplo, se puede borrar

% FIN DEL DOCUMENTO
\end{document}